% !TEX root = ../../cookbook.tex
\newpage
\recipe{Carbonara Pasta\label{carbonara}}{}{}{}
\vspace*{-0.5 cm}

\subsection*{Ingredients \underline{per person}}
\begin{ingredients}
	\item 80-120 g (2.8-4.2 oz) of pasta, according to the anticipated appetite of those at the table\footnote{Spaghetti are the traditional choice, but I personally prefer rigatoni. An excellent brand of pasta which is often found outside of Italy is \textit{Rummo}.}
	\item 2 egg yolks, plus 1 whole egg per 2 people\footnote{Different schools of thought disagree on the use of egg yolks, whole eggs, or a combination of the two. An excellent Carbonara can also be made using 1 whole egg per person, plus one egg ``\textit{for the pot}''. If favoring the yolk-based version, the remaining egg whites can be used to prepare ``\nameref{meringues}'' (p. \pageref{meringues}).}
	\item 40 g (1.4 oz) \textit{guanciale}, cut into short strips\footnote{\textit{Guanciale} is an Italian cured pork product made from a pig's cheek. If \textit{guanciale} is not available to you, you can substitute it with \textit{pancetta}, but know that this will change the ``soul'' of your carbonara. \textit{Pancetta} is made from pork belly and, among other processing differences, is not smoked like bacon, which is your last option to substitute \textit{guanciale}. In the US, \textit{pancetta} is typically available in many grocery stores.}
	\item 15 g (0.5 oz) grated Pecorino Romano\footnote{Pecorino Romano (literally, ``Roman sheep cheese'') is a hard, aged Italian cheese made from sheep milk. In the US, pecorino is often sold at Costco, but can substituted with Parmesan cheese. Just like with \textit{guanciale}, this substitution will inevitably change the fundamental nature of your final dish, offering a less salty and pungent result.}
	\item Salt and pepper
\end{ingredients}


\medskip

% --- Preparation ---
Cook the \textbf{pasta} just beyond al dente in salted water\footnote{As a rule, the salt level of the cooking water should correspond to the desired saltiness of the pasta cooked in it.}.
Fry the \textbf{guanciale} in a large skillet, which will later need to accommodate all the pasta.
Remove the guanciale from the skillet when golden but still soft, leaving all the rendered fat in the skillet, and turn off the heat.
Meanwhile, in a bowl, thoroughly mix the \textbf{egg yolks/eggs}, grated \textbf{Pecorino}, a little \textbf{salt}\footnote{The Pecorino and guanciale will provide plenty of natural saltiness on their own.}, and as much \textbf{pepper} as desired with a fork.
Once the pasta is cooked, reserve a ladleful (or a small glass) of the cooking water before draining, then transfer the pasta to the skillet with the guanciale fat.
Pour a little of the cooking water into the skillet with the pasta and, stirring continuously, add the egg and Pecorino mixture.
This is the stage that separates a good Carbonara from scrambled eggs: the egg should thicken into an emulsion together with the fat and water, ideally becoming pasteurized by the residual heat of the skillet, but it must absolutely not cook to the point of curdling\footnote{The ultimate goal of a creamy sauce is achieved by skillfully combining continuous stirring, the addition of more cooking water, and---only if the skillet has cooled to the point where the sauce no longer thickens---turning the heat back on to the very lowest setting.}.
Once the desired consistency has been achieved, add the guanciale and serve with additional pepper and grated Pecorino.

%\vspace*{0.75 cm}

% --- Description ---
\begin{recipeblock}
	
\end{recipeblock}